\scanpage{67}

Donc les \deleted{\unsure{poly}} fonctions \unsure{polygôniales|polynômiales} invariantes ont une base sur $k$ formée de $1$ et des

\[\begin{aligned}
&\deleted{\varphi_\alpha(x,y)}\\
\varphi_{\alpha,i}(x,y)&=x^\alpha y^{\alpha+in}+y^\alpha x^{\alpha+in}\\
&=x^\alpha y^\alpha(x^{in}+y^{in})\\
&=(xy)^\alpha\big(\underbrace{(x^n)^i+(y^n)^i}\big)
\end{aligned}\]

avec $\alpha\in\mathbb N$, $i\in\mathbb N^*$

\deleted{\unsure{Soit}} Or si $X,Y$ sont des indéterminées, le Th. des fonctions sym. nous dit que

\[\begin{aligned}
X^i+Y^i&=\underset{\inserted{\substack{\searrow\ \text{\unsure{polygône|polynôme} universel à coeff}\ \in\mathbb Z\\
\text{en }2\text{ variables }\sigma_1,\sigma_2}}}{S_i}(\sigma_1(X,Y),\sigma_2(X,Y))\\
&=S_i(X+Y,XY)
\end{aligned}\]

donc

\[(x^n)^i+(y^n)^i=S_i(x^n+y^n,(xy)^n)\]

et par suite

\[k[x,y]^{\overset{\inserted{\nearrow\,=\operatorname{Aut}(\Pi_{\zeta,k})}}{D_{\unsure{\zeta|G|{}}\deleted{\unsure{,k|\zeta}}}}}
=k[x,y]^{\overset{\inserted{\substack{\nearrow\,\text{groupe }\mu_n\mathbin{.}\mathbb Z/2\mathbb Z\ (\unsure{\tfrac12\,\text{direct}|\tfrac12\,\text{dirct}})\\
\text{opérant sur }k^2\text{ de façon évidente}}}}{\Delta_n}}=k[\sigma_2,S_n]\]

où

\[\left\{\begin{aligned}
&\sigma_2(x,y)=xy\\
&S_n(x,y)=x^n+y^n
\end{aligned}\right.\]

\uline{NB} Les calculs marchent pour tt $n\in\mathbb N^*$ (pas besoin que $n\geq3$) et sur tt anneau de base -- le meilleur de tous étant $\mathbb Z$ ! Si $n$ est infini, il faut revenir au cas d'un $\zeta$

% archive sheet number: 66
